\chapter{ไฟฟ้าสถิตและกลไกแรงดึงดูดทางไฟฟ้าของสีย้อมเซลล์จุลินทรีย์}

\begin{planbox}[สัปดาห์ที่ 8 — เวลาปฏิบัติการ 3 ชั่วโมง (สอดคล้องกับ CLO1, CLO3 และ CLO4)]
\textbf{1. วัตถุประสงค์การทดลอง}
\begin{itemize}
    \item ทดสอบแรงกระทำทางไฟฟ้าสถิตตามกฎของคูลอมบ์ (ประจุชนิดเดียวกันผลักกัน ประจุต่างชนิดกันดึงดูดกัน) ด้วยชุดอิเล็กโทรสโคป
    \item เปรียบเทียบประสิทธิภาพการยึดเกาะของสีย้อมประจุบวก (Basic Dye เช่น คริสตัลไวโอเลตและเมทิลีนบลู) กับสีย้อมประจุลบ (Acidic Dye เช่น นิกโรซินและอีโอซิน) บนพื้นผิวแบบจำลองผนังเซลล์ที่มีประจุสุทธิเป็นลบ ($\zeta\text{-potential} < 0$)
    \item อธิบายเหตุผลทางฟิสิกส์ของเทคนิคการย้อมสีแบบธรรมดา (Simple Staining) และการย้อมสีพื้นหลังแบบลบ (Negative Staining) ในงานจุลชีววิทยาได้อย่างถูกต้อง
\end{itemize}
\textbf{2. ความเชื่อมโยงกับผลลัพธ์การเรียนรู้ของรายวิชา (CLOs) และหลักสูตร (PLOs)}
\begin{itemize}
    \item \textbf{CLO1 (PLO2)}\quad อธิบายกฎของคูลอมบ์และธรรมชาติของประจุไฟฟ้าสุทธิบนผนังเซลล์แบคทีเรียได้
    \item \textbf{CLO3 (PLO1)}\quad ปฏิบัติงานกับสารละลายสีย้อมจุลชีววิทยาด้วยความระมัดระวัง สวมถุงมือยาง และกำจัดน้ำล้างสีย้อมลงในภาชนะเฉพาะตามหลักความปลอดภัยห้องปฏิบัติการ
    \item \textbf{CLO4 (PLO5, PLO6)}\quad สังเกต บันทึกภาพวาดจากกล้องจุลทรรศน์ และอภิปรายผลร่วมกันในกลุ่มได้อย่างเป็นระบบ
\end{itemize}
\end{planbox}

\section{หลักการทฤษฎีและการเชื่อมโยงสู่งานจุลชีววิทยา}

แรงดึงดูดและแรงผลักระหว่างประจุไฟฟ้า $q_1$ และ $q_2$ ที่อยู่ห่างกันเป็นระยะ $r$ อธิบายได้ด้วย\textbf{กฎของคูลอมบ์} (\textbf{Coulomb's Law})\index{กฎของคูลอมบ์ (Coulomb's Law)} โดยประจุไฟฟ้าชนิดเดียวกันจะออกแรงผลักกัน ส่วนประจุไฟฟ้าต่างชนิดกันจะออกแรงดึงดูดเข้าหากัน ในสารละลายน้ำที่มีค่าคงตัวไดอิเล็กทริกสัมพัทธ์สูง ($\varepsilon_r \approx 80$) ขนาดของแรงไฟฟ้าสถิตคำนวณได้จากสมการที่ \eqref{eq:lab07_coulomb}

\begin{equation}
    F_e = \frac{k\,|q_1 q_2|}{\varepsilon_r\,r^2}
    \label{eq:lab07_coulomb}
\end{equation}

ในสภาวะที่เป็นกลาง ($\text{pH} \approx 7.0$) ผนังเซลล์ของแบคทีเรียเกือบทุกชนิดมีหมู่ฟอสเฟต ($\text{PO}_4^{3-}$) และหมู่คาร์บอกซิล ($\text{COO}^-$) แตกตัวอยู่หนาแน่น ทำให้ผิวเซลล์แบคทีเรียมี\textbf{ประจุไฟฟ้าสุทธิเป็นลบ} ($\text{Zeta Potential} \approx -15\text{ ถึง }-45\text{ mV}$) เสมอ ดังนั้น
\begin{enumerate}
    \item \textbf{สีย้อมประจุบวก (Basic Dye / Cationic Dye)}\quad เช่น คริสตัลไวโอเลต ($\text{CV}^+$) เมทิลีนบลู ($\text{MB}^+$) และซาฟรานิน ($\text{Saf}^+$)\index{สีย้อมประจุบวก (Basic Dye)} จะถูกดึงดูดเข้ายึดเกาะกับผนังเซลล์แบคทีเรียที่มีประจุลบโดยตรง ทำให้เห็นตัวเซลล์ติดสีชัดเจนบนพื้นหลังใส
    \item \textbf{สีย้อมประจุลบ (Acidic Dye / Anionic Dye)}\quad เช่น นิกโรซิน (Nigrosin) และอีโอซิน ($\text{Eosin}^-$)\index{สีย้อมประจุลบ (Acidic Dye)} จะถูกผนังเซลล์แบคทีเรียผลักออกด้วยแรงคูลอมบ์ ทำให้สีไม่สามารถซึมเข้าสู่ตัวเซลล์ได้ แต่จะไปเคลือบอยู่รอบนอก เกิดเป็นภาพเซลล์ใสสว่างตัดกับพื้นหลังสีเข้ม เรียกว่า \textbf{การย้อมสีพื้นหลังแบบลบ} (\textbf{Negative Staining}) ซึ่งเหมาะสำหรับส่องดูแคปซูลของแบคทีเรีย
\end{enumerate}

\begin{conceptbox}[สรุปสูตรคำนวณฉบับเข้าใจง่ายสำหรับปฏิบัติการที่ 7]
\begin{itemize}
    \item \textbf{ขนาดของแรงคูลอมบ์ในอากาศ}\quad $F_0 = \dfrac{k\,|q_1 q_2|}{r^2}$ เมื่อ $k = 9.0 \times 10^9\text{ N}\cdot\text{m}^2/\text{C}^2$
    \item \textbf{ขนาดของแรงคูลอมบ์ในน้ำบริสุทธิ์ ($\varepsilon_r = 80$)}\quad $F_{\text{water}} = \dfrac{F_0}{80}$
    \item \textbf{กฎการจับคู่ประจุในการย้อมสีเซลล์ (ประจุผิวเซลล์เป็นลบเสมอที่ pH 7)}
    \begin{itemize}
        \item ผิวเซลล์ ($-$) $+$ สีย้อมประจุบวก ($+$) $\longrightarrow$ \textbf{ดึงดูดกัน} (ตัวเซลล์ติดสีเข้ม พื้นหลังใส)
        \item ผิวเซลล์ ($-$) $+$ สีย้อมประจุลบ ($-$) $\longrightarrow$ \textbf{ผลักกัน} (ตัวเซลล์ใสสว่าง พื้นหลังมืด)
    \end{itemize}
\end{itemize}
\end{conceptbox}

\begin{worked}[การเปรียบเทียบขนาดแรงดึงดูดทางไฟฟ้าในอากาศและในน้ำ]
\textbf{โจทย์} สมมติให้ไอออนบวกของสีย้อมคริสตัลไวโอเลต ($q_1 = +1.6 \times 10^{-19}\text{ C}$) อยู่ห่างจากหมู่ฟอสเฟตประจุลบของผนังเซลล์ ($q_2 = -1.6 \times 10^{-19}\text{ C}$) เป็นระยะทาง $r = 1.0\text{ nm}$ ($1.0 \times 10^{-9}\text{ m}$) หากแรงดึงดูดในสุญญากาศมีค่า $F_0 = 2.30 \times 10^{-10}\text{ N}$ จงคำนวณขนาดของแรงดึงดูดเมื่ออยู่ในตัวกลางที่เป็นน้ำ ($\varepsilon_r = 80$)

\textbf{วิธีทำทีละขั้นตอน}
\begin{enumerate}
    \item นำค่าแรงในสุญญากาศ $F_0 = 2.30 \times 10^{-10}\text{ N}$ หารด้วยค่าคงตัวไดอิเล็กทริกสัมพัทธ์ของน้ำ $\varepsilon_r = 80$
    \begin{equation*}
        F_{\text{water}} = \frac{2.30 \times 10^{-10}}{80} = 2.88 \times 10^{-12}\text{ N (หรือ } 2.88\text{ pN)}
    \end{equation*}
\end{enumerate}
\textbf{คำตอบ} ในตัวกลางน้ำ แรงดึงดูดทางไฟฟ้ามีค่าเท่ากับ $2.88 \times 10^{-12}\text{ N}$ ซึ่งยังมากพอที่จะยึดโมเลกุลสีย้อมประจุบวกให้ติดแน่นกับผนังเซลล์แบคทีเรีย
\end{worked}

\begin{figure}[H]
\centering
\resizebox{0.86\textwidth}{!}{%
\begin{tikzpicture}[>=Stealth, font=\small]
    % Left box: Basic Dye (+)
    \draw[fill=boxBlueBg, draw=physRoyalBlue, line width=1.1pt, rounded corners=6pt] (0,0) rectangle (6.6,4.2);
    \node[font=\bfseries\color{physRoyalBlue}] at (3.3,3.8) {(ก) การย้อมด้วยสีย้อมประจุบวก (Basic Dye)};
    \draw[fill=purple!45, draw=purple!80!black, line width=1.4pt, rounded corners=12pt] (1.6,1.4) rectangle (5.0,2.7);
    \node[font=\bfseries\color{white}] at (3.3,2.05) {ตัวเซลล์ติดสีม่วงเข้ม};
    \foreach \x in {1.9, 2.6, 3.3, 4.0, 4.7} {
        \node[font=\scriptsize\bfseries\color{crimsonAlert}] at (\x,2.88) {$-$};
        \node[font=\scriptsize\bfseries\color{crimsonAlert}] at (\x,1.22) {$-$};
    }
    \node[circle, fill=physRoyalBlue, text=white, inner sep=1.2pt, font=\scriptsize\bfseries] (cv1) at (1.1,3.2) {$\text{CV}^+$};
    \node[circle, fill=physRoyalBlue, text=white, inner sep=1.2pt, font=\scriptsize\bfseries] (cv2) at (5.5,3.2) {$\text{CV}^+$};
    \draw[->, line width=1.1pt, bioEmerald!85!black] (cv1) -- (1.75,2.75);
    \draw[->, line width=1.1pt, bioEmerald!85!black] (cv2) -- (4.85,2.75);
    \node[font=\footnotesize\bfseries\color{bioEmerald!85!black}] at (3.3,0.45) {แรงคูลอมบ์ดึงดูด ($-$ ดึงดูด $+$)};

    % Right box: Acidic Dye (-)
    \draw[fill=charcoal!85, draw=charcoal, line width=1.1pt, rounded corners=6pt] (7.4,0) rectangle (14.0,4.2);
    \node[font=\bfseries\color{physAmber}] at (10.7,3.8) {(ข) การย้อมสีพื้นหลังแบบลบ (Negative Stain)};
    \draw[fill=white, draw=charcoal!60, line width=1.4pt, rounded corners=12pt] (9.0,1.4) rectangle (12.4,2.7);
    \node[font=\bfseries\color{charcoal}] at (10.7,2.05) {ตัวเซลล์ใสสว่าง};
    \foreach \x in {9.3, 10.0, 10.7, 11.4, 12.1} {
        \node[font=\scriptsize\bfseries\color{crimsonAlert}] at (\x,2.88) {$-$};
        \node[font=\scriptsize\bfseries\color{crimsonAlert}] at (\x,1.22) {$-$};
    }
    \node[circle, fill=crimsonAlert, text=white, inner sep=1.2pt, font=\scriptsize\bfseries] (ng1) at (8.3,2.8) {$\text{Nig}^-$};
    \node[circle, fill=crimsonAlert, text=white, inner sep=1.2pt, font=\scriptsize\bfseries] (ng2) at (13.1,2.8) {$\text{Nig}^-$};
    \draw[->, line width=1.1pt, physAmber] (8.85,2.6) -- (8.1,3.25);
    \draw[->, line width=1.1pt, physAmber] (12.55,2.6) -- (13.3,3.25);
    \node[font=\footnotesize\bfseries\color{physAmber}] at (10.7,0.45) {แรงคูลอมบ์ผลักออก ($-$ ผลัก $-$)};
\end{tikzpicture}%
}
\caption{เปรียบเทียบกลไกแรงไฟฟ้าสถิตระหว่างผิวเซลล์แบคทีเรียประจุลบกับ (ก) สีย้อมประจุบวก และ (ข) สีย้อมประจุลบ}
\label{fig:lab07_setup}
\end{figure}

\section{วัสดุอุปกรณ์และขั้นตอนการปฏิบัติการ}

\begin{equipbox}
\begin{enumerate}
    \item ชุดอิเล็กโทรสโคปแผ่นโลหะพร้อมแท่งพีวีซี แท่งอะคริลิกใส และผ้าสักหลาด จำนวน $1$ ชุด
    \item แผ่นกระดาษกรองเซลลูโลส (จำลองพื้นผิวประจุลบ) และสไลด์แก้วสะอาด จำนวน $4$ แผ่น
    \item สารละลายสีย้อมประจุบวก เมทิลีนบลู ($0.1\%$ Methylene Blue) และคริสตัลไวโอเลต จำนวนอย่างละ $1$ ขวดหยด
    \item สารละลายสีย้อมประจุลบ อีโอซิน ($1\%$ Eosin Y) และหมึกอินเดีย/นิกโรซิน จำนวนอย่างละ $1$ ขวดหยด
    \item สไลด์ตัวอย่างเซลล์ยีสต์ \textit{Saccharomyces cerevisiae} และกล้องจุลทรรศน์แบบใช้แสง จำนวน $1$ ชุด
\end{enumerate}
\end{equipbox}

\begin{warningbox}
สวมถุงมือยางตลอดเวลาที่หยดสารละลายสีย้อมเพื่อป้องกันสีย้อมประจุบวกย้อมติดผิวหนัง และเทน้ำล้างสไลด์ลงในขวดรวบรวมของเสียสีย้อมที่จัดเตรียมไว้ ห้ามเทลงท่อน้ำทิ้งโดยตรง
\end{warningbox}

\subsection{ขั้นตอนการทดลอง}
\begin{enumerate}
    \item \textbf{ตอนที่ 1 การทดสอบแรงกระทำระหว่างประจุไฟฟ้าด้วยอิเล็กโทรสโคป}\quad นำแท่งพีวีซีถูกับผ้าสักหลาดแล้วจ่อใกล้จานโลหะของอิเล็กโทรสโคป สังเกตการกางออกของแผ่นโลหะบาง จากนั้นทดสอบแรงผลักและแรงดึงดูดระหว่างแท่งพีวีซีกับแท่งอะคริลิกที่ถูด้วยผ้าไหม บันทึกผลใน\mbox{ตารางที่} 7.1
    \item \textbf{ตอนที่ 2 การเปรียบเทียบการยึดเกาะของสีย้อมประจุบวกและประจุลบบนพื้นผิวประจุลบ}\quad หยดสีย้อมเมทิลีนบลู (ประจุบวก) และอีโอซิน (ประจุลบ) ลงบนแผ่นเซลลูโลสที่มีประจุลบในปริมาณเท่ากัน ($20\ \mu\text{L}$) ทิ้งไว้ $2$ นาที แล้วล้างด้วยน้ำกลั่นไหลผ่านเบาๆ นาน $30$ วินาที สังเกตความเข้มของสีที่หลงเหลืออยู่
    \item \textbf{ตอนที่ 3 การส่องเปรียบเทียบ Simple Staining และ Negative Staining ของเซลล์ยีสต์}\quad หยดสารแขวนลอยยีสต์ลงบนสไลด์ 2 แผ่น แผ่นที่ 1 ย้อมด้วยเมทิลีนบลู (ประจุบวก) และแผ่นที่ 2 ผสมกับนิกโรซิน (ประจุลบ) แล้วไถสไลด์ให้บาง นำไปส่องด้วยกล้องจุลทรรศน์กำลังขยาย $400\times$ บันทึกลักษณะสีของตัวเซลล์และพื้นหลังลงใน\mbox{ตารางที่} 7.2
\end{enumerate}

\section{ใบบันทึกผลการทดลองและคำถามท้ายปฏิบัติการ}

\begin{worksheetbox}[ใบบันทึกผลการทดลองที่ 7\quad ไฟฟ้าสถิตและกลไกของสีย้อมเซลล์จุลินทรีย์]
\textbf{กลุ่มปฏิบัติการที่} \dotfill \textbf{วันที่ทำการทดลอง} \dotfill\\[4pt]
\textbf{รายชื่อสมาชิกกลุ่ม} 1. \dotfill 2. \dotfill 3. \dotfill

\vspace{4pt}
\textbf{ตารางที่ 7.1\quad ผลการทดสอบแรงกระทำทางไฟฟ้าสถิตและการชะล้างสีย้อมบนพื้นผิวประจุลบ}
\begin{center}
\footnotesize\setlength{\tabcolsep}{3.5pt}
\begin{tabularx}{\linewidth}{>{\raggedright\arraybackslash}p{4.2cm} c >{\raggedright\arraybackslash}X >{\raggedright\arraybackslash}X}
\toprule
\textbf{ชุดการทดลอง} & \textbf{ชนิดประจุ} & \textbf{ผลที่สังเกตได้ (ดึงดูด/ผลัก/การติดสี)} & \textbf{คำอธิบายตามกฎของคูลอมบ์} \\
\midrule
1. แท่งพีวีซี 2 แท่งที่ถูผ้าสักหลาดแล้วนำมาเข้าใกล้กัน & $(-)$ กับ $(-)$ & \dotfill & \dotfill \\[6pt]
2. แท่งพีวีซีถูผ้าสักหลาด เข้าใกล้แท่งอะคริลิกถูผ้าไหม & $(-)$ กับ $(+)$ & \dotfill & \dotfill \\[6pt]
3. หยด Methylene Blue บนแผ่นประจุลบแล้วล้างน้ำ & ผิว $(-)$ สี $(+)$ & \dotfill & \dotfill \\[6pt]
4. หยด Eosin Y บนแผ่นประจุลบแล้วล้างน้ำกลั่น & ผิว $(-)$ สี $(-)$ & \dotfill & \dotfill \\
\bottomrule
\end{tabularx}
\end{center}

\vspace{4pt}
\textbf{ตารางที่ 7.2\quad ผลการสังเกตเซลล์ยีสต์ภายใต้กล้องจุลทรรศน์ ($400\times$) ด้วยสีย้อมต่างชนิดประจุ}
\begin{center}
\footnotesize\setlength{\tabcolsep}{3.5pt}
\begin{tabularx}{\linewidth}{>{\raggedright\arraybackslash}p{3.8cm} c >{\raggedright\arraybackslash}X >{\raggedright\arraybackslash}X >{\raggedright\arraybackslash}X}
\toprule
\textbf{ชนิดของสีย้อมที่ใช้} & \textbf{ประจุของไอออนสี} & \textbf{สีของตัวเซลล์ยีสต์} & \textbf{สีของพื้นหลังสไลด์} & \textbf{ประเภทการย้อม (Simple / Negative)} \\
\midrule
เมทิลีนบลู (Methylene Blue) & ประจุบวก $(+)$ & \dotfill & \dotfill & \dotfill \\[8pt]
นิกโรซิน (Nigrosin) & ประจุลบ $(-)$ & \dotfill & \dotfill & \dotfill \\
\bottomrule
\end{tabularx}
\end{center}

\vspace{4pt}
\textbf{สรุปและอภิปรายผลการทดลอง}\\[4pt]
\dotfill\\[6pt]
\dotfill
\end{worksheetbox}

\subsection{คำถามท้ายปฏิบัติการ}
\begin{enumerate}
    \item เพราะเหตุใดผนังเซลล์ของแบคทีเรียและยีสต์ในสภาวะค่า $\text{pH}$ เป็นกลาง ($\text{pH} \approx 7.0$) จึงมีประจุไฟฟ้าสุทธิเป็นลบ และประจุลบนี้เกิดจากหมู่ฟังก์ชันทางเคมีใดเป็นหลัก
    \item หากนักศึกษาต้องการตรวจสอบว่าแบคทีเรียสายพันธุ์หนึ่งสร้างแคปซูล (Capsule) หุ้มเซลล์หรือไม่ เหตุใดจึงนิยมใช้เทคนิคการย้อมสีพื้นหลังแบบลบด้วยสีนิกโรซินควบคู่กับการย้อมสีตัวเซลล์ แทนการย้อมสีแบบธรรมดาเพียงอย่างเดียว
\end{enumerate}
